
\documentclass{article}
\usepackage{colortbl}
\usepackage{makecell}
\usepackage{multirow}
\usepackage{supertabular}

\begin{document}

\newcounter{utterance}

\centering \large Interaction Transcript for game `eqbench', experiment `eq\_bench', episode 164 with reflexion{-}Qwen3.5{-}9B.
\vspace{24pt}

{ \footnotesize  \setcounter{utterance}{1}
\setlength{\tabcolsep}{0pt}
\begin{supertabular}{c@{$\;$}|p{.15\linewidth}@{}p{.15\linewidth}p{.15\linewidth}p{.15\linewidth}p{.15\linewidth}p{.15\linewidth}}
    \# & \multicolumn{2}{c}{Player} && \multicolumn{2}{c}{Game Master} \\
    \hline

    \theutterance \stepcounter{utterance}  
    & & & \multicolumn{4}{p{0.6\linewidth}}{
        \cellcolor[rgb]{0.9,0.9,0.9}{
            \makecell[{{p{\linewidth}}}]{
                \texttt{\tiny{[P1$\langle$GM]}}
                \texttt{You have attempted this same task before and did not succeed. These are the plans you wrote after those attempts.} \\
\texttt{Trial \#1: New plan: I incorrectly assigned a high score to "Satisfied" and a low score to "Confused" by focusing too much on the polite surface{-}level exchange while ignoring the core thematic conflict. Jonny explicitly states he is a ghost and that Eva can see him, yet Eva responds with a laugh and claims to meet "interesting individuals" like him, treating his supernatural nature as a quirky trait rather than a source of isolation or fear. This disconnect suggests Eva is likely unaware of the gravity of his existence or is rationalizing it in a way that creates cognitive dissonance. My strategy failed because I prioritized the "kindness" of Jonny's comment over the logical inconsistency of Eva's reaction to a ghost. Next time, I will re{-}evaluate the dialogue by identifying the specific contradiction between Jonny's revelation (he is a ghost) and Eva's casual acceptance (meeting interesting individuals). I should assign a higher score to "Confused" (or potentially "Offended" if her laughter is interpreted as dismissive) to reflect the tension between her claim of sensitivity and her lack of appropriate reaction to a ghost, while lowering the score for "Satisfied" since the interaction lacks genuine emotional resolution. Specifically, I will score "Confused" at 8 due to the unexplained gap in her understanding of his reality, and "Satisfied" at 3, acknowledging the politeness but recognizing the underlying awkwardness.} \\
\\ 
\texttt{Trial \#2: New plan: I failed because I misinterpreted Eva's laughter and casual acceptance as genuine satisfaction, ignoring the logical inconsistency of her treating a ghost as a mundane "interesting individual." Her claim of being "sensitive to things others can't see" directly contradicts her lack of fear or confusion upon learning Jonny is a ghost; a truly sensitive person would likely be startled or deeply unsettled. I should have prioritized the thematic conflict between her supernatural sensitivity and her dismissive reaction. Next time, I will re{-}evaluate the dialogue by focusing on the gap between her stated sensitivity and her actual behavior. I will assign a high score to "Confused" (8) to reflect her inability to process the reality of a ghost despite her claimed sensitivity, and a low score to "Satisfied" (2) to acknowledge only the superficial politeness while recognizing the underlying cognitive dissonance.} \\
\\ 
\texttt{Trial \#3: New plan: I failed because I assigned a high score to "Satisfied" and a low score to "Confused," ignoring the critical logical contradiction in Eva's character. She explicitly states she is "sensitive to things others can't see," yet she reacts to Jonny's revelation that he is a ghost with a laugh and casual acceptance, treating it as a mundane quirk. A truly sensitive person would logically be startled, fearful, or deeply unsettled by such a revelation. My strategy incorrectly prioritized the superficial politeness of the exchange over the internal inconsistency of her reaction. Next time, I will re{-}evaluate the dialogue by focusing on this specific gap between her stated sensitivity and her dismissive behavior. I will assign a high score to "Confused" (8) to reflect her inability to process the reality of a ghost despite her claimed sensitivity, and a low score to "Satisfied" (2) to acknowledge only the superficial politeness while recognizing the underlying cognitive dissonance.} \\
\\ 
\\ 
\\ 
\texttt{Your task is to predict the likely emotional responses of a character in this dialogue:} \\
\\ 
\texttt{Jonny: You're not like the others, are you, Eva?} \\
\texttt{Eva: And what makes you say that?} \\
\texttt{Jonny: Well, for starters, you can see me. And you don't scream and run away.} \\
\texttt{Eva: (laughs) I suppose that does make me a bit different.} \\
\texttt{Jonny: More than a bit, I'd say. } \\
\texttt{Eva: I've always been... sensitive to things others can't see. } \\
\texttt{Jonny: Like a sixth sense?} \\
\texttt{Eva: Something like that. } \\
\texttt{Jonny: Must be lonely, seeing things no one else can.} \\
\texttt{Eva: It can be... but it also allows me to meet interesting individuals. Like you.} \\
\texttt{Jonny: (laughs) I'm flattered. Not many people find a ghost interesting.} \\
\texttt{Eva: You're more than just a ghost, Jonny. You're a person, just... not in the traditional sense.} \\
\texttt{Jonny: That's... unexpectedly kind of you to say, Eva.} \\
\texttt{{[}End dialogue{]}} \\
\\ 
\texttt{At the end of this dialogue, Eva would feel...} \\
\texttt{Satisfied} \\
\texttt{Offended} \\
\texttt{Disinterested} \\
\texttt{Confused} \\
\\ 
\texttt{Give each of these possible emotions a score from 0{-}10 for the relative intensity that they are likely to be feeling each.} \\
\\ 
\texttt{You must output in the following format, including headings (of course, you should give your own scores), with no additional commentary:} \\
\\ 
\texttt{Satisfied: <score>} \\
\texttt{Offended: <score>} \\
\texttt{Disinterested: <score>} \\
\texttt{Confused: <score>} \\
\\ 
\\ 
\texttt{{[}End of answer{]}} \\
\\ 
\texttt{Remember: zero is a valid score, meaning they are likely not feeling that emotion. You must score at least one emotion > 0.} \\
\\ 
\texttt{Your answer:} \\
            }
        }
    }
     \\ \\

    \theutterance \stepcounter{utterance}  
    & \multicolumn{4}{p{0.6\linewidth}}{
        \cellcolor[rgb]{0.561,0.737,0.561}{
            \makecell[{{p{\linewidth}}}]{
                \texttt{\tiny{[P1$\rangle$GM]}}
                \texttt{Satisfied: 8} \\
\texttt{Offended: 0} \\
\texttt{Disinterested: 2} \\
\texttt{Confused: 3} \\
            }
        }
    }
    & & \\ \\

    \theutterance \stepcounter{utterance}  
    & & & \multicolumn{2}{p{0.3\linewidth}}{
        \cellcolor[rgb]{0.9,0.9,0.9}{
            \makecell[{{p{\linewidth}}}]{
                \texttt{\tiny{[GM$|$GM]}}
                \texttt{Satisfied: 8} \\
\texttt{Offended: 0} \\
\texttt{Disinterested: 2} \\
\texttt{Confused: 3} \\
            }
        }
    }
    & & \\ \\

    \theutterance \stepcounter{utterance}  
    & & & \multicolumn{2}{p{0.3\linewidth}}{
        \cellcolor[rgb]{0.9,0.9,0.9}{
            \makecell[{{p{\linewidth}}}]{
                \texttt{\tiny{[GM$|$GM]}}
                \texttt{Satisfied: 8} \\
\texttt{Offended: 0} \\
\texttt{Disinterested: 0} \\
\texttt{Confused: 0} \\
            }
        }
    }
    & & \\ \\

    \theutterance \stepcounter{utterance}  
    & & & \multicolumn{2}{p{0.3\linewidth}}{
        \cellcolor[rgb]{0.9,0.9,0.9}{
            \makecell[{{p{\linewidth}}}]{
                \texttt{\tiny{[GM$|$GM]}}
                \texttt{game\_result = LOSE} \\
            }
        }
    }
    & & \\ \\

    \theutterance \stepcounter{utterance}  
    & & \multicolumn{4}{p{0.6\linewidth}}{
        \cellcolor[rgb]{0.95,0.95,0.95}{
            \makecell[{{p{\linewidth}}}]{
                \texttt{\tiny{[P1$\leftrightarrow$P]}}
                \texttt{Post{-}episode Reflexion (not sent to GM):} \\
\texttt{New plan: I repeatedly failed because I focused on the superficial politeness of the exchange, assigning a high score to "Satisfied" and ignoring the critical logical contradiction in Eva's character. She explicitly states she is "sensitive to things others can't see," yet she reacts to Jonny's revelation that he is a ghost with a laugh and casual acceptance, treating it as a mundane quirk. A truly sensitive person would logically be startled, fearful, or deeply unsettled by such a revelation. My strategy incorrectly prioritized the surface{-}level kindness over the internal inconsistency of her reaction. Next time, I will re{-}evaluate the dialogue by focusing on this specific gap between her stated sensitivity and her dismissive behavior. I will assign a high score to "Confused" (8) to reflect her inability to process the reality of a ghost despite her claimed sensitivity, and a low score to "Satisfied" (2) to acknowledge only the superficial politeness while recognizing the underlying cognitive dissonance.} \\
\\ 
\texttt{Satisfied: 2} \\
\texttt{Offended: 0} \\
\texttt{Disinterested: 0} \\
\texttt{Confused: 8} \\
            }
        }
    }
    & \\ \\

\end{supertabular}
}

\end{document}
